\setcounter{ex}{0}
\setcounter{paragraph}{0}
\PhanI
%\loai{Tính cơ năng khi biết phương trình dao động}
\begin{ex} %[3P1B5-2][Loai1]  
Một chất điểm khối lượng $ m = 150$ g, dao động điều hoà với phương trình $x = 3\cos(2t + \dfrac{\pi}{3})\ cm$. Cơ năng trong dao động điều hoà của chất điểm này là
\choice
{0,27 J}
{3,2 J}
{0,32 J}
{\True 0,27 mJ}
\loigiai{
Ta có: $W=\dfrac{m\omega^2A^2}{2} = \dfrac{0,15.2^2.0,03^2}{2} = 2,7.10^{-4}\ J=0,27\ mJ$.
}%<MyLT1>
%\haidong{2}
\end{ex} \haidong{20}

%\loai{Tính động năng cực đại}
\begin{ex} %[3P1B5-2][Loai1] 
Một vật nhỏ có khối lượng bằng 100 g, dao động điều hòa với biên độ bằng 5 cm và chu kì bằng $\pi$ s. Động năng cực đại của vật là
\choice
{$6.10^{-4}\ J$}
{$7\ J$}
{$4\ J$}
{\True $5.10^{-4}\ J$}
\loigiai{
Ta có: $\omega = \dfrac{2\pi}{T} = 2\ rad/s\Rightarrow W_{\text{đmax}} = \dfrac{m\omega^2A^2}{2} = \dfrac{0,1.2^2.0,05^2}{2} = 5.10^{-4}\ J.$
}
%\haidong{2}
\end{ex} \haidong{20}

%\loai{Từ cơ năng tìm các đại lượng}
\begin{ex} %[3P1B5-2][Loai1] 
Một vật dao động điều hoà có phương trình $x = 4\cos(3t-\dfrac{\pi}{6})$ cm. Cơ năng của vật là $7,2.10^{-3}$ J. Khối lượng của vật là
\choice
{\True 1 kg}
{2 kg}
{0,1 kg}
{0,2 kg}
\loigiai{
Ta có: $W=\dfrac{1}{2}kA^2=\dfrac{1}{2}m\omega^2.A^2 \Rightarrow 7,2.10^{-3} = \dfrac{1}{2}m.3^2.0,04^2\Rightarrow m=1\ kg$.
}
%\haidong{2}
\end{ex} \haidong{20}

\begin{ex}%[3P1B5-2][Loai1]
Một con lắc lò xo dao động điều hòa. Biết lò xo có độ cứng 36 N/m và vật nhỏ có khối lượng 100 g. Lấy $\pi^2 = 10$. Động năng của con lắc biến thiên theo thời gian với tần số là
\choice
{3 Hz}
{\True 6 Hz}
{1 Hz}
{12 Hz}
\loigiai{
\begin{itemize}
\item Ta có: $\omega = \sqrt{\dfrac{k}{m}} = \sqrt{\dfrac{36}{0,1}} = 6\pi \ rad/s\Rightarrow f = \dfrac{\omega}{2\pi} = 3\ Hz$.
\item Động năng của con lắc biến thiên theo thời gian với tần số $2f=6$ Hz.
\end{itemize}
}
%\haidong{2}
\end{ex} \haidong{20}
%\loai{Từ động năng, thế năng tìm các đại lượng}
\begin{ex}%[3P1K5-2][Loai1] 
Một vật $m = 200$ g dao động điều hoà. Trong khoảng thời gian một chu kì vật đi được một đoạn 40 cm. Tại vị trí $x = 5$ cm thì động năng của vật là 0,375 J. Chu kì dao động là
\choice
{$t=0,045$ s}
{$t=0,02$ s}
{\True $t=0,28$ s}
{$t=0,14$ s}
\loigiai{
\begin{itemize}
\item Ta có: $4A = 40\ cm\Rightarrow A = 10\ cm$.
\item Tại vị trí có $x=5\ cm=A/2$ thì $W_\text{đ}=3W_t$ $\Rightarrow W_t = \dfrac{W_{\text{đ}}}{3} = \dfrac{1}{2}kx^2 = 0,125\ J \Rightarrow k = 100\ N/m \Rightarrow T = 0,28\ s$.
\end{itemize}
}
%\haidong{4}
\end{ex} \haidong{20}


%\loai{Liên hệ năng lượng ở các vị trí khác nhau}
\begin{ex} %[3P1B5-2][Loai1]  
Một con lắc lò xo mà lò xo nhẹ có độ cứng 100 N/m và vật nhỏ dao động điều hòa. Khi vật có động năng 0,01 J thì nó cách vị trí cân bằng 1 cm. Khi nó có động năng 0,005 J thì nó cách vị trí cân bằng là
\choice
{6 cm}
{4,5 cm}
{\True $\sqrt{2}$ cm}
{3 cm}
\loigiai{Hỏi khi vật có động năng 0,005 J
\begin{align*}
W=W_{\text{đ}}+\dfrac{kx^2}{2}\Rightarrow \begin{cases}
W=0,01+\dfrac{100.0,01^2}{2}.\\ 
W=0,005+\dfrac{100.x_2^2}{2}. 
\end{cases}\Rightarrow x_2=0,01\sqrt{2} m.
\end{align*}}
%\haidong{4}
\end{ex} \haidong{20}

%\loai{Biểu thức thế năng}
\begin{ex} %[3P1K5-2][Loai1] 
Một con lắc lò xo có $k = 30$ N/m, đang dao động điều hòa theo phương trình $x = 4\cos(3t + \dfrac{\pi}{6})$ cm. Thế năng của chất điểm dao động theo phương trình nào sau đây?
\choice
{$W_t = 12\cos(3t + \dfrac{\pi}{6}) + 12$ mJ}
{\True $W_t = 12\cos(6t + \dfrac{\pi}{3}) + 12$ mJ}
{$W_t = -12\cos(6t + \dfrac{\pi}{3}) + 12$ mJ}
{$W_t = 24\cos(6t + \dfrac{\pi}{3})$ mJ}
\loigiai{
\begin{itemize}
\item Ta có: $W=\dfrac{kA^2}{2} = \dfrac{30.0,04^2}{2} = 0,024\ J=24\ mJ$.
\item Mà $W_t = \dfrac{W}{2}.\cos(6t + \dfrac{\pi}{3}) + \dfrac{W}{2}\Rightarrow W = 12\cos(6t + \dfrac{\pi}{3}) + 12\ mJ$.
\end{itemize}
}
%\haidong{4}
\end{ex} \haidong{20}

%\loai{Biểu thức động năng}
\begin{ex}%[3P1K5-2][Loai1]
Một chất điểm khối lượng $ m = 200$ g, dao động điều hòa theo phương trình $x = 5\cos(0,5\pi t + \dfrac{2\pi}{3})$ cm. Động năng của chất điểm biến thiên theo phương trình nào sau đây?
\choice
{$W_{\text{đ}} = -0,617\cos(0,5\pi t + \dfrac{2\pi}{3}) + 0,617$ mJ}
{$W_{\text{đ}} = -0,617\cos(\pi t + \dfrac{\pi}{3}) + 0,617$ mJ}
{\True $W_{\text{đ}} = -0,308\cos(\pi t + \dfrac{4\pi}{3}) + 0,308$ mJ}
{$W_{\text{đ}} = 0,617\cos(\pi t + \dfrac{\pi}{3}) + 0,617$ mJ}
\loigiai{
\begin{itemize}
\item Ta có: $W=\dfrac{m\omega^2A^2}{2} = \dfrac{0,2.(0,5\pi)^2.0,05^2}{2} = 0,617.10^{-3}\ J=0,617\ mJ$.
\item Mà $W_{\text{đ}} = -\dfrac{W}{2}\cos(\pi t + \dfrac{4\pi}{3}) + \dfrac{W}{2} \Rightarrow W_{\text{đ}} = -0,308\cos(\pi t + \dfrac{4\pi}{3}) + 0,308$ mJ.
\end{itemize}
}
%\haidong{4}
\end{ex} \haidong{20}
%\loai{Tổng hợp}



\begin{ex}%[3P1K8-1][Loai1]
Một vật có khối lượng 1 kg dao động điều hoà dọc theo trục Ox với phương trình: $x = 10\cos(\omega t + \pi)\ cm$. Thời gian ngắn nhất vật đi từ vị trí $x = - 5\ cm$ đến vị trí $x = + 5\ cm$ là $\dfrac{\pi}{30}$ s. Cơ năng dao động của vật bằng
\choice
{\True $0,5\ J$}
{$5\ J$}
{$0.3\ J$}
{$3\ J$}
\loigiai{
\immini[thm]{\begin{itemize}
\item Ta biểu diễn vị trí vật có li độ $-5\ cm$ và $+5\ cm$ trên đường tròn lượng giác như hình vẽ.
\item Thời gian ngắn nhất sẽ ứng với vật đi từ $M_1$ đến $M_2\Rightarrow \Delta \varphi = \dfrac{\pi}{6}+\dfrac{\pi}{6} = \dfrac{\pi}{3};\ \Delta t = \dfrac{T}{6} = \dfrac{\pi}{30} \Rightarrow T = \dfrac{\pi}{5} \Rightarrow \omega = 10\ rad/s
\Rightarrow W = \dfrac{m\omega^2A^2}{2} = \dfrac{1.10^2.0,1^2}{2} = 0,5\ J$.
\end{itemize}}{\begin{tikzpicture}[scale=1,thick,>=stealth']
\tkzInit[ymin=-3,ymax=3,xmin=-3,xmax=3]\tkzClip
\tkzDefPoints{-2.5/0/m, 2.5/0/n, 0/0/o, 0/2.5/p, 0/-2.5/q}
\tkzDefShiftPoint[o](-120:2){1}
\tkzDefShiftPoint[o](-60:2){2}
\tkzDefPointBy[projection=onto m--n](1) \tkzGetPoint{h}
\tkzDefPointBy[projection=onto m--n](2) \tkzGetPoint{h'}
\draw(o)circle(2cm);
\tkzDrawSegments[dashed](1,h 2,h')
\tkzDrawSegments(p,q)
\tkzDrawSegments[->](m,n)
\tkzDrawSegments[blue,->](o,1 o,2)
\tkzDrawPoints[size=3](o,1,h,2,h')
\tkzMarkAngles[size=0.7cm,arrows=->](1,o,2)
\node at (1)[below left]{$M_1$};
\node at (2)[below right]{$M_2$};
\node at (2,0)[above right]{$10$};
\node at (h)[above]{$-5$};
\node at (h')[above]{$5$};
\end{tikzpicture}}
}%<MyLT1>
%\haidong{5}
\end{ex} \haidong{20}

\begin{ex}%[3P1-19.1K][Loai1]
\immini[thm]{Cho một con lắc lò xo gồm một vật nhỏ khối lượng $ m = 100 $ g gắn với lò xo, đang dao động điều hòa. Sự phụ thuộc của li độ vào thời gian được biểu diễn trên đồ thị như hình vẽ. Lấy $\pi ^2=10$. Động năng của vật tại vị trí có li độ $ x = 2 $ cm là
\choice
{0,36 J}
{\True 0,48 J}
{0,52 J}
{0,5 J}}{\begin{tikzpicture}[very thick,>=stealth']
\pgfplotsset{width=7cm,height=5cm,compat=1.4}
\begin{axis}[
axis lines=middle, xmin=0, xmax=5, xstep=1,
ymin=-4.5, ymax=5, ystep=0.1,
grid=major,%Vẽ chế độ lưới theo trục tọa độ Oxy,
x grid style={dashed, black},
axis line style={very thick,Mapcolor}, % <--- dòng thêm vào
xlabel=$t~(s)$,
xlabel style={above},
xtick ={1,2,3,4},
xticklabels={},
y grid style={dashed, black},
extra y ticks={0},
extra y tick label={0},
ylabel=$x~(cm)$,
ylabel style={above},
ytick={-4,-2,0,2, 4},
yticklabels={$-10$,,,,10},]
\addplot[line width=1.2pt, domain=0:4, samples=400, Mapcolor]{4*cos(deg(0.5*pi*x-pi/2))};
\addplot[Mapcolor,mark=*,mark options={fill=white,mark size=1.5pt}, nodes near coords={0,2}]
coordinates {(4, 0)};
\end{axis}
\end{tikzpicture}}
\loigiai{
\begin{itemize}
\item Từ đồ thị ta thấy: $\dfrac{T}{2}=0,1\ s\Rightarrow T=0,2\ s$.
\item Độ cứng của lò xo: $ k = 100 $ N/m.
\item Áp dụng định luật bảo toàn năng lượng: $W = W_t + W_{\text{đ}}\Leftrightarrow \dfrac{kA^2}{2} = \dfrac{kx^2}{2} + W_{\text{đ}} \Rightarrow W_{\text{đ}} = 0,48\ J$.
\end{itemize}
}
%\haidong{4}
\end{ex} \haidong{20}
\begin{ex} %[3P1K5-2][Loai1]   
Một chất điểm đang dao động điều hòa. Khi vừa qua khỏi vị trí cân bằng một đoạn S động năng của chất điểm là 0,091 J. Đi tiếp một đoạn 2S thì động năng chỉ còn 0,019 J và nếu đi thêm một đoạn S ($A>3S$) nữa thì động năng bây giờ là
\choice
{0,042 J}
{0,096 J}
{\True 0,036 J}
{0,032 J}
\loigiai{
\begin{itemize}
\item Ta có: $\begin{cases} 0,091=W-\dfrac{1}{2}kS^2\\
0,019=W-\dfrac{1}{2}k(3S)^2\end{cases}\Rightarrow \begin{cases}W=0,1\ J=\dfrac{1}{2}kA^2\\
\dfrac{1}{2}kS^2=0,009\ J\end{cases}\Rightarrow \dfrac{A}{S}=\dfrac{10}{3}$. 
\item Khi vật đi thêm đoạn S nữa là 4S thì vật đã đi qua biên và quay lại một đoạn $\dfrac{2}{3}S$ nên vật ở li độ $x=\dfrac{8}{3}S$: \begin{align*}W_{\text{đ}}=W-\dfrac{1}{2}k\left(\dfrac{8}{3}S\right)^2=W-\dfrac{64}{9}.\dfrac{1}{1}kS^2=0,036\ J\end{align*}

\end{itemize}
}
%\haidong{4}
\end{ex} \haidong{20}

\begin{ex} %[3P1B5-3][Loai1] 
Một chất điểm dao động dọc theo trục Ox. Phương trình dao động là $x = 2\sin10t$ cm. Li độ x của chất điểm khi động năng bằng 3 lần thế năng có độ lớn bằng
\choice
{2 cm}
{$\sqrt{2}$ cm}
{\True 1 cm}
{0,707 cm}
\loigiai{
Ta có: $W_{\text{đ}} = 3W_t \Rightarrow W_t = \dfrac{1}{4} W\Rightarrow x^2 = \dfrac{1}{4}A^2 \Rightarrow |x| = \dfrac{A}{2} = 1\ cm$.
}

%\haidong{4}
\end{ex} \haidong{20}
\begin{ex}%[3P1K5-3][Loai1]
Một vật có khối lượng $m = 200$ g gắn với một lò xo có độ cứng $k = 20$ N/cm. Từ vị trí cân bằng kéo vật đến li độ 5 cm rồi truyền cho nó tốc độ 5 m/s. Sau đó, vật dao động điều hòa. Vị trí của vật mà tại đó động năng bằng 3 lần thế năng cách vị trí cân bằng là
\choice
{1 cm}
{\True $2,5\sqrt{2}$ cm}
{3 cm }
{4 cm}
\loigiai{
\begin{itemize}
\item Ta có: $k = 20$ N/cm $= 2000$ N/m $\Rightarrow$ $\omega  = 100$ rad/s.
\item Mặt khác:  $A=\sqrt{x^2+\dfrac{v^2}{\omega^2}}=5\sqrt{2}$ cm.
\item Khi $W_{\text{đ}}=3W_t\Rightarrow \left| x \right|=\dfrac{A}{2}=2,5\sqrt{2}$ cm.
\end{itemize}
}
%\haidong{4}
\end{ex} \haidong{20}




%\loai{Cho li độ, tìm tỉ số động năng và thế năng}


\begin{ex}%[3P1K5-3][Loai1]
Một vật nhỏ khối lượng 200 g đang dao động điều hòa với chu kì là 0,2 s và cơ năng toàn phần bằng 0,25 J (mốc thế năng tại vị trí cân bằng). Lấy $\pi^2 = 10$. Tại li độ 2,5 cm, tỉ số động năng và thế năng là
\choice
{1}
{4}
{\True 3}
{2}
\loigiai{
Ta có: 
\begin{align*}
&\omega = \dfrac{2\pi}{T} = 10\pi \ rad/s\Rightarrow W_{t} = \dfrac{m\omega^2x^2}{2} = \dfrac{0,2.(10\pi)^2.0,025^2}{2} = 0,0625\ J\\
\Rightarrow\ &W_{\text{đ}} = 0,25-0,0625 = 0,1875\ J\Rightarrow \dfrac{W_{\text{đ}}}{W_t} = \dfrac{0,1875}{0,0625} = 3.
\end{align*}
}
%\haidong{4}
\end{ex} \haidong{20}

%\loai{Tỉ số cơ năng với động năng hoặc thế năng}



%\loai{Cho vận tốc, gia tốc tìm liên hệ cơ năng với động năng và thế năng}
\begin{ex}%[3P1K5-3][Loai1]
Vật nhỏ của một con lắc lò xo dao động điều hòa theo phương ngang, mốc thế năng tại vị trí cân bằng. Khi gia tốc của vật có độ lớn bằng một nửa độ lớn gia tốc cực đại thì tỉ số giữa động năng và thế năng của vật là
\choice
{$\dfrac{1}{2}$}
{\True 3}
{2}
{$\dfrac{1}{3}$}
\loigiai{
Ta có: $\left| a \right|=\dfrac{a_{\max}}{2}=\dfrac{\omega^2A}{2}\Rightarrow \left| x \right|=\dfrac{A}{2}\Rightarrow W_{\text{đ}}=3W_t$.
}
%\haidong{7}
\end{ex} \haidong{20}


%\loai{Từ liên hệ giữa năng lượng tìm liên hệ giữa x và v}


%\loai{Từ liên hệ năng lượng tìm các đại lượng khác x và v}


\begin{ex}%[3P1K5-3][Loai1]
Một chất điểm dao động điều hòa trên trục Ox, trong một phút thực hiện được 150 dao động toàn phần. Khi vật có tọa độ 2 cm thì tốc độ là $10\pi$  cm/s. Tại thời điểm $t=0,$ vật có động năng bằng thế năng, sau đó vật có li độ tăng và động năng tăng. Phương trình dao động của vật là
\choice
{$x=4\cos  \left( 2\pi t+\dfrac{\pi}{4} \right)\ cm$}
{$x=2\sqrt{2}\cos  \left( 5\pi t+\dfrac{\pi}{4} \right)\ cm$}
{$x=2\sqrt{2}\cos  \left( 2\pi t-\dfrac{3\pi}{4} \right)\ cm$}
{\True $x=2\sqrt{2}\cos  \left( 5\pi t-\dfrac{3\pi}{4} \right)\ cm$}
\loigiai{
\begin{itemize}
\item $\begin{cases}
T =  \dfrac{60}{150}=0,4\ s\Rightarrow \omega =5\pi\ rad/s.\\ 
A=\sqrt{x^2+\dfrac{v^2}{\omega^2}}=2\sqrt{2}\ cm.
\end{cases}$ 
\item Tại $t = 0:\ W_{\text{đ}}=W_t\Leftrightarrow x=\pm \dfrac{A\sqrt{2}}{2}$.
\item Li độ tăng $\Rightarrow x=\pm \dfrac{A\sqrt{2}}{2}\oplus$.
\item Động năng tăng $\Rightarrow x=-\dfrac{A\sqrt{2}}{2}\oplus \Rightarrow \varphi =-\dfrac{3\pi}{4}$.
\item Vậy: $x=2\sqrt{2}\cos  \left( 5\pi t-\dfrac{3\pi}{4} \right)$ cm.
\end{itemize}
}
%\haidong{8}
\end{ex} \haidong{20}

















%\loai{Đồ thị năng lượng hình sin}

%\loai{Thời điểm đầu tiên liên quan động năng - thế năng của dao động điều hoà}


%\loai{Thời điểm $W_{\ct{đ}}=nW_t$ lần thứ n}


%\loai{Khoảng thời gian lặp lại liên quan đến năng lượng}




%\loai{Khoảng thời gian trong một chu kì thoả mãn hệ thức năng lượng}



%\loai{Khoảng thời gian liên quan năng lượng}



%\loai{Năng lượng của vật tại trạng thái sau, trước}
\begin{ex}%[3P1K9-2][Loai1]
Nếu vào thời điểm ban đầu, vật dao động điều hòa đi qua vị trí cân bằng thì vào thời điểm $\dfrac{T}{12}$, tỉ số giữa động năng và thế năng của dao động là
\choice
{1}
{\True 3}
{2}
{$\dfrac{1}{3}$}
\loigiai{
\begin{itemize}
\item Giả sử pha ban đầu của vật là $\dfrac{\pi}{2}$.
\item Sau khoảng thời gian $t=\dfrac{T}{12}$ thì vật quay được góc $\Delta \varphi = \dfrac{\pi}{6}$.
\item Thời điểm $\dfrac{T}{12}$ vật ở vị trí góc $\dfrac{2\pi}{3}$ trên VTLG $\Rightarrow x = -\dfrac{A}{2}\Rightarrow W_t = \dfrac{1}{4}W \Rightarrow W_\text{đ} = \dfrac{3}{4}W\Rightarrow \dfrac{W_\text{đ}}{W_t} = 3$.
\end{itemize}
}
%\haidong{5}
\end{ex} \haidong{20}




%\loai{Đồ thị liên hệ thế năng và động năng}

\begin{ex}%[3P1-20.1K][Loai1]
\immini[thm]{Con lắc lò xo có vật nhỏ có khối lượng 400 g đang dao động điều hòa. Hình bên là đồ thị biểu diễn sự phụ thuộc của động năng $W_{\text{đ}}$ của con lắc theo thời gian t. Tại $t = 0$, vật đang chuyển động theo chiều dương. Lấy $\pi^2 = 10$. Phương trình dao động của vật là
\choice
{$ x=5\cos  \left( 2\pi t+\dfrac{\pi}{3} \right) $ cm}
{$ x=10\cos  \left( \pi t+\dfrac{\pi}{6} \right) $ cm}
{\True $ x=5\cos  \left( 2\pi t-\dfrac{\pi}{3} \right) $ cm}
{$ x=10\cos  \left( \pi t-\dfrac{\pi}{3} \right) $ cm}}{\begin{tikzpicture}[very thick,>=stealth']
\pgfplotsset{width=7cm,height=5cm,compat=1.4}
\begin{axis}[
axis lines=middle, xmin=0, xmax=5, xstep=1,
ymin=-4.5, ymax=5, ystep=0.1,
grid=major,%Vẽ chế độ lưới theo trục tọa độ Oxy,
x grid style={dashed, black},
axis line style={very thick,Mapcolor}, % <--- dòng thêm vào
xlabel=,
xlabel style={below},
axis x line shift=4,
xtick ={1/3,2/3,3/3,4/3,5/3,6/3,7/3,8/3,9/3,10/3,11/3,12/3,13/3,14/3},
xticklabels={,,,$\dfrac{1}{6}$,,,,,,,,,,t(m/s)},
y grid style={dashed, black},
extra y ticks={-4},
extra y tick label={0},
ylabel=$W_{\text{đ}} (mJ)$,
ylabel style={above},
ytick={-4,-2,0,2, 4},
yticklabels={,,10,,20},]
\addplot[line width=1.2pt, domain=0:4.5, samples=400, Mapcolor]{4*cos(deg(0.5*pi*x+2*pi/6))};
\end{axis}
\end{tikzpicture}}

\loigiai{
\begin{itemize}
\item Từ đồ thị:  $\dfrac{T'}{3}=\dfrac{1}{6}\Rightarrow T'=0,5\ s \Rightarrow T=2T'=1\ s \Leftrightarrow \omega  = 2\pi\ rad/s$.  
\item $W_{t\max}=W=0,02\ J=\dfrac{1}{2}m\omega^2A^2\Rightarrow A = 5\ cm$.  
\item Tại $t = 0$: $W_{\text{đ}} = 3W_t \Rightarrow  \left| x \right|=\dfrac{A}{2}$ và đi theo chiều dương.
\item Lại thấy sau $t = 0$, $W_{\text{đ}}$ giảm $\Rightarrow  x=\dfrac{A}{2}$ và đi theo chiều dương $ \Rightarrow \varphi =-\dfrac{\pi}{3}$. 
\item Vậy: $ x=5\cos  \left( 2\pi t-\dfrac{\pi}{3} \right)\ cm$.
\end{itemize}
}
%\haidong{5}
\end{ex} \haidong{20}


  \setcounter{ex}{0}
\PhanII
\begin{ex}%book
Một chất điểm dao động điều hòa trên trục $Ox$ với biên độ $A$.
\choiceTF[t]
{\True Cứ mỗi chu kì dao động của vật, có bốn thời điểm thế năng bằng động năng}
{Động năng bằng thế năng khi vật ở vị trí có li độ $x=\pm \dfrac{A}{2}$}
{Động năng của vật tăng và thế năng giảm khi vật đi từ vị trí cân bằng đến vị trí biên}
{Động năng giảm, thế năng tăng khi vật đi từ vị trí biên đến vị trí cân bằng}
\loigiai{
\begin{itemchoice}
\itemch
\itemch
\itemch
\itemch
\end{itemchoice}

}
\end{ex} \haidong{20}

\begin{ex}%book
Xét dao động điều hoà có phương trình $x=A\cos(\omega t)\ cm$.
\choiceTF[t]
{Đồ thị biểu diễn thế năng theo vận tốc là một đoạn thẳng đi qua gốc tọa độ}
{Thế năng của chất điểm tại thời điểm $t$ có biểu thức là \linebreak $W_{t}=\dfrac{W_0}{2}(1-\cos 2 \omega t)$}
{\True Khi động năng của chất điểm giảm thì thế năng của nó tăng}
{Động năng biến thiên vuông pha so với thế năng}
\loigiai{
\begin{itemchoice}
\itemch
\itemch
\itemch
\itemch
\end{itemchoice}

}
\end{ex} \haidong{20}
\begin{ex}%book
Xét dao động điều hoà có phương trình $x=A\cos(\omega t)\ cm$.
\choiceTF[t]
{Động năng vuông pha với cơ năng của vật}
{\True Trong một chu kì, có hai thời điểm động năng cực đại}
{Trong một chu kì, có hai thời điểm động năng gấp ba lần thế năng}
{\True Thế năng gấp ba lần động năng tại vị trí vật có li độ $x=\pm \dfrac{A \sqrt{3}}{2}$}
\loigiai{
\begin{itemchoice}
\itemch
\itemch
\itemch
\itemch
\end{itemchoice}

}
\end{ex} \haidong{20}
\begin{ex}%book
\immini[thm]{Một vật dao động điều hoà có các đồ thị biểu diễn mối quan hệ giữa cơ năng $W$, động năng $W_{\text{đ}}$, thế năng $W_{t}$ có dạng là một trong các đường ở hình bên.}{\begin{tikzpicture}[draw=Mapcolor,scale=.7,thick,>=stealth']
\tkzDefPoints{0/0/o, 4/0/x, 0/4/y}
\tkzDrawSegments[->](o,x o,y)
\tkzLabelPoint[below](x){$x$}
\tkzLabelPoint[above](y){$y$}
\tkzLabelPoint[below left](o){$O$}
\draw[Mapcolor,very thick] (0,0) -- (3,3);
\draw[Mapcolor,very thick](0,2.5) -- (3.5,2.5);
\draw[Mapcolor,very thick] (2,3) -- (2,0);
\draw[very thick](0,3) -- (3,0);
\node at (3.25,3.35) {$I$};
\node at (4,2.5) {$II$};
\node at (2,3.35) {$III$};
\node at (0.5,2) {$IV$};
\end{tikzpicture}}
\choiceTF[t]
{\True Đồ thị biểu diễn mối quan hệ giữa động năng và thế năng là đường IV}
{Đồ thị biểu diễn mối quan hệ giữa động năng và bình phương li độ là đường III với trục $O y$ biểu diễn động năng}
{Đồ thị biểu diễn mối quan hệ giữa động năng và li độ là đường I với trục $Ox$ biểu diễn cơ năng}
{Đồ thị biểu diễn mối quan hệ giữa thế năng và tốc độ là đường II với trục $O y$ biểu diễn tốc độ}
\loigiai{
\begin{itemchoice}
\itemch
\itemch
\itemch
\itemch
\end{itemchoice}
}
\end{ex} \haidong{20}

\setcounter{ex}{0}
\PhanIII
\begin{ex}
Động năng của một chất điểm dao động điều hoà được mô tả bởi phương trình sau: $W_{\text{đ}}=0,8 \sin ^2\left(6 \pi t+\dfrac{\pi}{6}\right)\ J$ (t tính bằng s). Thế năng của chất điểm tại thời điểm $t=1 \ s$ là bao nhiêu jun?
\shortans[oly]{0,6}
\loigiai{
0,6
}
%\haidong{6}
\end{ex} \haidong{20}
\begin{ex}
Một vật nhỏ thực hiện dao động điều hòa theo phương trình $x=10 \sin \left(4 \pi t+\dfrac{\pi}{2}\right)\ cm$ với $t$ tính bằng giây. Động năng của vật đó biến thiên với chu kì bằng bao nhiêu giây?
\shortans[oly]{0,25}
\loigiai{
0,25
}
%\haidong{6}
\end{ex} \haidong{20}
\begin{ex}
Một vật nhỏ khối lượng $100 \ g$ dao động theo phương trình $x=8 \cos 10 t\ cm$. Động năng cực đại của vật bằng bao nhiêu $m J$?
\shortans[oly]{32}
\loigiai{
32
}
%\haidong{6}
\end{ex} \haidong{20}
\begin{ex}
Một vật nhỏ dao động điều hòa dọc theo trục $O x$. Khi vật cách vị trí cân bằng một đoạn $2 \ cm$ thì động năng của vật là $0,48 \ J$. Khi vật cách vị trí cân bằng một đoạn $6 \ cm$ thì động năng của vật là $0,32 \ J$. Biên độ dao động của vật bằng bao nhiêu $cm$?
\shortans[oly]{10}
\loigiai{
10
}
%\haidong{6}
\end{ex} \haidong{20}
\begin{ex} %[3P1-20.1K][Loai1]
\immini[thm]{Một con lắc lò xo có vật nhỏ khối lượng 0,1 kg dao động điều hòa trên trục Ox với phương trình $x = A\cos\omega t$ (cm). Đồ thị biểu diễn động năng theo bình phương li độ như hình vẽ. Lấy $\pi^2 = 10$. Tốc độ trung bình của vật trong 1 chu kì là bao nhiêu cm/s?
}{\begin{tikzpicture}[draw=Mapcolor,scale=1,thick,>=stealth']
\tkzDefPoints{0/0/o, 2.6/0/x, 0/2.5/i, 0/2/a, 2/0/b}
\tkzDrawSegments[->](o,x o,i)
\tkzDrawSegments[Mapcolor,very thick](a,b)
\node at (a)[left]{$0,08$};
\node at (b)[below]{$16$};
\node at (x)[above]{$x^2\ (cm^2)$};
\node at (i)[right]{$W_{\text{đ}}\ (J)$};
\node at (o)[below left]{$O$};
\end{tikzpicture}
}
\shortans[oly]{80}
\loigiai{
\begin{itemize}
\item Từ đồ thị ta thấy $x_{max}^2 = 16 = A^2 \Rightarrow A = 4\ cm$.
\item $W_{\text{đmax}} = W = \dfrac{1}{2}m\omega^2A^2 = 0,08\ J \Rightarrow \omega = 10\pi\ rad/s \Rightarrow T = 0,2\ s$.
\item Vậy $\bar{v}=\dfrac{4A}{T} = 80\ cm/s$.
\end{itemize}
}
%\haidong{5}
\end{ex} \haidong{20}
\begin{ex} %[3P1-20.1K][Loai1] 
\immini[thm]{\nguon{MH - 2017} Hình vẽ sau là đồ thị biểu diễn sự phụ thuộc của thế năng đàn hồi $W_{\text{đh}}$ của một con lắc lò xo vào thời gian t. Tần số dao động của con lắc bằng bao nhiêu Hz?
}{\begin{tikzpicture}[very thick,>=stealth']
\pgfplotsset{width=7cm,height=5cm,compat=1.4}
\begin{axis}[
axis lines=middle, xmin=0, xmax=5, xstep=1,
ymin=-4.5, ymax=5, ystep=0.1,
grid=major,%Vẽ chế độ lưới theo trục tọa độ Oxy,
x grid style={dashed, black},
axis line style={very thick,Mapcolor}, % <--- dòng thêm vào
xlabel=$t~(ms)$,
xlabel style={above},
axis x line shift=4,
xtick ={1/2,1,1.5,2,2.5,3,3.5,4,4.5},
xticklabels={,5,,10,,15,,20},
y grid style={dashed, black},
extra y ticks={0},
extra y tick label={0},
ylabel=$W_{dn}$,
ylabel style={above},
ytick={-4,-2,0,2, 4},
yticklabels={,,,},]
\addplot[line width=1.2pt, domain=0:4.5, samples=400, Mapcolor]{4*cos(deg(0.5*pi*x-pi/2))};
\end{axis}
\end{tikzpicture}}
\shortans[oly]{25}
\loigiai{
\begin{itemize}
\item Nhìn vào đồ thị ta có chu kì dao động của thế năng đàn hồi là $T' = 20\ ms$.
\item Chu kì dao động của con lắc lò xo: $T = 40\ ms$.
\item Tần số dao động của con lắc lò xo: $f = \dfrac{1}{T} = 25\ Hz$.
\end{itemize}
}
%\haidong{9}
\end{ex} \haidong{20}
